\subsection{The conormal calculation and ordinary-triple equality}

\begin{statement}{P4.30}{Proved}
Let $C$ be a compact irreducible reduced curve in the smooth
threefold $Y$ of Part~4.2. Suppose its normalization is an elliptic
curve $E$, and every point of $C$ is unibranch. Write $P$ for its
ideal in $Y$ and $P^{(2)}$ for the ideal of functions vanishing to
order at least two at the generic point of $C$. Then
\[
 \chi(\mathcal O_Y/P^{(2)})\le0.
\]
Equality forces $C$ to be smooth. The assertion uses
$c_1(TY)\cdot C=0$, which follows from $b_2(Y)=0$.
\end{statement}
\begin{proof}
Let $j:E\to Y$ be the normalization map followed by inclusion.
The derivative map
\[
 j^*\Omega_Y\longrightarrow\Omega_E
\]
has image $\Omega_E(-R)$ for an effective divisor $R$.
Its kernel $K$ is locally free of rank two: dividing the derivative
row by its common zero gives a surjection onto a line bundle.
Since $\deg j^*\Omega_Y=0$ and $\deg\Omega_E=0$, we have
$\deg K=\deg R$.

Differentiation induces an injection
\[
 Q:=P/P^{(2)}\longrightarrow j_*K.
\]
Indeed its kernel consists precisely of functions with zero first
normal derivative at a generic smooth point of $C$, which is the
defining order-two condition for $P^{(2)}$. Both sides agree at the
generic point. Their quotient has finite length, denoted by $\tau$.
Riemann--Roch on the smooth elliptic curve gives
\[
 \chi(Q)=\deg R-\tau.
\]
We prove $\tau\ge\deg R$ by local equations, including at a
singular point of $C$.

Parametrize its sole branch by $t$. Choose ambient coordinates
vanishing at the point, and let
$m=\min_i\operatorname{ord}_t x_i(t)$ be the multiplicity.
The common order of the three derivative entries is $m-1$,
so the coefficient of $R$ is $m-1$. The saturated kernel $K$
is a direct summand of the free rank-three module over
$\mathbb C\{t\}$. Consequently a generator of $\bigwedge^2K$
has at least one unit Pl\"ucker coordinate. If $L$ is the
$\mathbb C\{t\}$-span of the differentiated defining equations,
then the colength of $L$ in $K$ is the minimum order of their
two-by-two minors. This is the diagonal elementary-divisor
calculation for a rank-two submodule over a discrete valuation
ring. The local length of $j_*K/Q$ is at least this colength,
since $Q\subset L\subset K$.

At a smooth point $m=1$ and there is nothing to prove.
At a singular point of embedding dimension three, every defining
equation belongs to the square of the ambient maximal ideal.
Every entry of its gradient therefore has order at least $m$
on the parametrization, and every two-by-two minor has order
at least $2m$. Hence the local quotient length is at least
$2m\ge m-1$.

At a singular point of embedding dimension two, an invertible
ambient coordinate change writes the ideal as $(z,h(x,y))$,
where $h$ is a reduced irreducible plane equation without a
linear term. The conormal image has a free $dz$ summand.
Its other generator is $dh$; both its plane coefficients have
order at least $m$. The primitive rank-one kernel of the plane
derivative row has a unit coefficient, so $dh$ is its generator
multiplied by a function of order at least $m$. The quotient
length is therefore at least $m\ge m-1$. Embedding dimension
one is the smooth case. These cases prove the local bound and
therefore $\chi(Q)\le0$.

The normalization exact sequence gives
$\chi(\mathcal O_C)=-\sum_p\delta(C,p)$.
Finally the exact sequence
\[
 0\longrightarrow P/P^{(2)}\longrightarrow
 \mathcal O_Y/P^{(2)}\longrightarrow\mathcal O_C
 \longrightarrow0
\]
gives
\[
 \chi(\mathcal O_Y/P^{(2)})
 =-\sum_p\delta(C,p)+\deg R-\tau\le0.
\]
Equality forces all delta invariants to vanish, so the
normalization is an isomorphism and $C$ is smooth.
\end{proof}

\begin{statement}{P4.31}{Proved}
Let $S$ be an irreducible singular support in Part~4.2 whose
normalization has ruled resolution, and suppose $e(S)=0$.
In the double-horizontal-divisor case, when the generic
transverse type is the ordinary triple point $xy(x-y)=0$,
the normalization is a smooth minimal elliptic ruled surface.
Its conductor image is a smooth elliptic curve, and its
reduced conductor is an unramified cover of that curve of
degree three. Every germ along the image is analytically
the product of an ordinary plane triple point with a disc.
\end{statement}
\begin{proof}
Use the minimal resolution $T\to N\to S$ and the section $H$
of P4R.7 and P4R.9. Put $g=g(H)$ and
$C=\operatorname{Sing}(S)_{\rm red}$.
The same fundamental-group and Euler argument as in P4.29
shows that $C$ is connected, has just one normalized
component, that this normalization $E$ has genus one, and
that its point-identification graph has no cycles. Thus
every point of $C$ is unibranch. All rational components
of the resolved pole support are contracted, and the
only strict conductor component is $H$. Its finite map
to $C$ has degree three, since an ordinary transverse
triple point has three normalization branches.
In particular $g\ge1$ by Riemann--Hurwitz for $H\to E$.

All exceptional curves of $T\to N$ are rational and
vertical for the ruling. Exactly the normalized-graph
descent and Leray calculation of P4.29 give a morphism
$N\to H$ and
\[
 R^1\rho_{N*}\mathcal O_T=0,\qquad
 \chi(\mathcal O_N)=\chi(\mathcal O_T)=1-g.
\]
Thus the normalization singularities are rational; no
smoothness assertion is used here.

Let $I$ be the conductor ideal, and let $C_{\rm sch}$
and $D_{\rm sch}$ be its zero schemes in $S$ and $N$.
Finite duality and the depth argument of P4.29 show that
$\mathcal O_{C_{\rm sch}}$ is Cohen--Macaulay of pure
dimension one. If $P\subset\mathcal O_Y$ is the reduced
ideal of $C$, the inverse image $J$ of $I$ in
$\mathcal O_Y$ therefore has no embedded associated
prime. Its only associated prime locally is $P$.
The generic ordinary triple point has conductor equal
to the square of its maximal ideal in the transverse
plane. To check this directly, normalize its three
lines. Constant functions on them have a common value;
their linear terms are the evaluations of one linear
form on three distinct directions, a two-dimensional
space. Quadratic forms evaluate surjectively on those
three directions. All terms of degree at least two
can therefore be chosen independently on the three
branches, and a nonzero constant or linear term cannot
belong to the conductor. This proves the stated square
ideal, including its scheme structure.

Consequently $J_P=P_P^2$, and in fact
\[
 J=P^{(2)}.
\]
For completeness, if $a\in P^{(2)}$, some $s\notin P$
satisfies $sa\in J$. Multiplication by $s$ is injective
on $\mathcal O_Y/J$, because its only associated prime
is $P$, so $a\in J$. The reverse inclusion follows by
localizing at $P$. Thus the conductor image is the
actual symbolic double curve to which P4.30 applies.

Put $L=\nu^*\omega_S$, which is torsion. Finite duality
gives $I=\omega_N\otimes L^{-1}$. Serre duality, the
rational-singularity comparison and Riemann--Roch on
$T$ give
\[
 \chi(I)=\chi(L)=\chi(\rho_N^*L)
   =\chi(\mathcal O_T)=1-g.
\]
Hence $\chi(\mathcal O_{D_{\rm sch}})=0$.
Taking Euler characteristics in the exact conductor
square, and using $\chi(\mathcal O_S)=0$, yields
\[
 \chi(\mathcal O_{C_{\rm sch}})=g-1.
\]
P4.30 now gives $g-1\le0$. Since $g\ge1$, equality
holds, $g=1$, and $C$ is smooth. Riemann--Hurwitz then
makes the degree-three map $H\to C$ unramified.

We finish by proving normalization smoothness from
the actual local morphisms. Over a point of $C$ the
three distinct points of $H$ remain distinct in $N$:
every exceptional block is vertical for the ruling
$T\to H$ and therefore cannot join different points
of $H$. The finite normalization consequently gives
three local irreducible branches of $S$. Each branch
is a hypersurface in $Y$ and is smooth at the generic
point of $C$. Its singularities are isolated, so the
hypersurface Cohen--Macaulay property and regularity
in codimension one make it normal. The corresponding
component of $N$ is that branch itself.

Write its three local equations as $f_1,f_2,f_3$.
The intersection of any two branches is supported on
$C$ and has generic multiplicity one, because the
generic transverse triple point is ordinary. On a
normal branch the other equation is a nonzero divisor.
Its quotient is a one-dimensional Cohen--Macaulay
ring, generically reduced with support $C$, so it is
reduced and is exactly $\mathcal O_C$. Thus $C$ is a
smooth Cartier curve in each branch. A local ring of
dimension two with a nonzero divisor whose quotient
is a regular one-dimensional local ring has embedding
dimension at most two: lift its one maximal-ideal
generator and add the nonzero divisor. It is regular.
All three branches are therefore smooth.

Any two branches meet transversely along $C$.
Choose their equations as coordinates $x,y$, with
$z$ a coordinate on $C$. The third equation is
$a(x,y,z)x+b(x,y,z)y$, where both $a$ and $b$ are
units by the pairwise transversality just proved.
The invertible coordinate change $X=ax$, $Y=-by$
writes the reduced union exactly as
$XY(X-Y)=0$. This is the asserted product triple
point, and its normalization is smooth. Away from
the conductor, normalization singularities were
already excluded by P4R.6. Finally P4.22 applies to
this smooth elliptic ruled normalization and proves
its minimality.
\end{proof}
